\section{関連研究}
転移強化学習では，ソース環境で得た知識をターゲット環境へ
再利用することで学習効率を改善する試みが広く研究されている．
なかでも過去に学習した方策を学習中の探索に組み込む「方策再利用」は，
環境構造が類似する場合に学習初期の性能を押し上げる一方，
不適切な再利用は負の転移を招くため，再利用の制御が重要となる．

Fern\'andez と Veloso は，
複数の過去方策を保持した方策ライブラリから，
ターゲット環境における学習を支援するために
過去方策を確率的に再利用する枠組みを提案した 
\cite{Fernandez2006}．
方策ライブラリ $\mathcal{L}=\{\Pi_1,\Pi_2,\ldots,\Pi_n\}$ 
と学習中の新方策 $\Pi_{\Omega}$ を候補とし，
エピソードごとに1つの方策を選択して実行する．
各方策 $\Pi_i$ には有効性指標 $W_i$（Reuse Gain）を保持し，
得られたリターンに基づいて更新する．
文献\cite{Fernandez2006}では，最大の $W$ を持つ方策指標として
\begin{equation}
j^\ast = \arg\max_{i=1,\ldots,n} W_i
\tag{3}
\end{equation}
を定義し，候補方策 $\Pi_j$ の選択確率を
\begin{equation}
P(\Pi_j)=\frac{\exp(\tau W_j)}{\sum_{p=0}^{n}\exp(\tau W_p)}
\tag{4}
\end{equation}
で与える（$p=0$ は新方策に対応）．

また，同研究では，選択した方策に常に従うのではなく，
再利用率 $\psi$ により探索を混合する実行戦略（$\pi$-reuse）を導入している．
具体的には，時刻 $t$ の状態 $s_t$ における行動 $a_t$ を
\begin{equation}
a_t =
\begin{cases}
\displaystyle \arg\max_{a} Q_{\Pi_j}(s_t,a)
& (\text{w/ prob. } \psi)\\[3pt]
\displaystyle \arg\max_{a} Q_{\Pi_{\mathrm{new}}}(s_t,a)
& (\text{w/ prob. } 1-\psi)
\end{cases}
\tag{5}
\end{equation}

のように決定し，
過去方策の過信を抑えつつターゲット環境での学習を進める．

以上を踏まえ，本研究はPRQ-Learningの枠組みを保ったまま
方策ライブラリからの方策選択部を拡張し，
ボルツマン選択に限定しない比較を行う点に特徴がある．
